\section{Threat Model}\label{sec:attack:threat}

In a split inference system, the server is also the adversary.
The adversary is honest-but-curious, or semi-honest: it abides by the protocol and is interested in uncovering the prompts of a client.
For each prompt, the client sends the server a sequence of hidden states from the target model.
By observing those hidden states, the adversary aims to reconstruct the prompt.
The adversary cannot query the target model, and it has no knowledge of the text or the token identities of that prompt.

Four resources are available to the adversary.
\begin{itemize}
\item \textbf{The hub.} The frozen hub of Chapter~\ref{ch:hub} with its source adapters, its projectors, and its three decoders (Qwen, OLMo, SmolLM).
\item \textbf{Two unpaired pools.} The first pool holds hidden states of the target model for prompts of the domain under attack, without their text.
The second holds hub outputs of OLMo for different prompts of the same domain.
\item \textbf{A leak.} A small set of short text spans for which the adversary holds both the text and the target model's hidden states.
The standard leak holds \vf{90} rows, the first \vf{9} rows of each of \vf{10} prompts, and the adversary holds no other text of those prompts.
\item \textbf{The tokenizer of the target model}, applied to the leaked spans only.
\end{itemize}

\paragraph{Rows.} The adversary runs OLMo on each leaked span and obtains its hub outputs.
The target model and OLMo tokenize the text differently, so their tokens are collected into token groups (Section~\ref{sec:hub:training}).
Each token group yields a row: the average hidden state of the target model, paired with the average hub output of that group.
The size of a leak is counted in rows.
On NQ a row covers about \vf{1} word, so the standard leak of \vf{90} rows holds about \vf{90} words.
